% =====================================================================
%  papers/paper2/paper2.tex
%
%  Paper 2 --- Selective inference for weights-matrix selection.
%
%  SCAFFOLD created by T-005.  The section skeleton, the labelled section
%  scheme (AC-012), the glossary wiring, the bibliography wiring and the
%  figure path are load-bearing and should be preserved.  The prose inside
%  each section is placeholder and belongs to the writing tasks.
%
%  Build:  cd papers && make paper2
%          GNU make is not installed on this machine; the equivalent and
%          actually-exercised driver is:  python papers/build.py paper2
%
%  AUTHORING RULE (D-005): edit this file with the Write/Edit tools only.
%  Bash heredocs collapse \\ to \ in this environment and corrupt tabulars.
% =====================================================================

\documentclass[11pt,a4paper]{article}

% =====================================================================
%  papers/shared/preamble.tex
%  Shared preamble for paper1 and paper2.  Owned by T-005 (scaffold).
%
%  Package set is restricted to the verified-working list:
%    amsmath, amssymb, amsthm, booktabs, graphicx, hyperref, natbib, geometry
%  Do not add packages without compiling to confirm they resolve.
%
%  AUTHORING RULE (D-005): this file and every other .tex in this tree are
%  written with the Write/Edit tools.  Bash heredocs collapse \\ to \ in this
%  environment and silently corrupt tabulars.
% =====================================================================

\usepackage[a4paper,margin=1.1in]{geometry}

\usepackage{amsmath}
\usepackage{amssymb}
\usepackage{amsthm}
\usepackage{booktabs}
\usepackage{graphicx}
\usepackage[authoryear,round]{natbib}

% hyperref last (except for its own configuration below).
\usepackage[
  colorlinks=true,
  linkcolor=black,
  citecolor=black,
  urlcolor=black,
  bookmarksnumbered=true,
  bookmarksopen=true,
  pdfstartview=FitH
]{hyperref}

% Figures live next to each paper's main .tex, so a bare relative path works
% for both documents when latexmk runs inside the paper directory.
\graphicspath{{figures/}{./figures/}}

\bibliographystyle{plainnat}
\setcitestyle{authoryear,round}

% ---------------------------------------------------------------------
%  Theorem-like environments
% ---------------------------------------------------------------------
\theoremstyle{plain}
\newtheorem{theorem}{Theorem}[section]
\newtheorem{proposition}[theorem]{Proposition}
\newtheorem{lemma}[theorem]{Lemma}
\newtheorem{corollary}[theorem]{Corollary}

\theoremstyle{definition}
\newtheorem{definition}[theorem]{Definition}
\newtheorem{example}[theorem]{Example}

\theoremstyle{remark}
\newtheorem{remark}[theorem]{Remark}

% ---------------------------------------------------------------------
%  \begin{keyinsight} ... \end{keyinsight}
%
%  Marks a contribution or a load-bearing idea so a reader skimming for the
%  new result can find it.  Built from base LaTeX primitives only -- no
%  framing package is in the verified set.
%  Optional argument overrides the label, e.g. \begin{keyinsight}[What is new]
% ---------------------------------------------------------------------
\newenvironment{keyinsight}[1][Key insight]{%
  \par\medskip
  \begin{list}{}{%
    \setlength{\leftmargin}{1.5em}%
    \setlength{\rightmargin}{1.5em}%
    \setlength{\listparindent}{0pt}%
    \setlength{\parsep}{0pt}%
    \setlength{\topsep}{0pt}%
    \setlength{\partopsep}{0pt}%
  }%
  \item[]%
  \noindent\rule{\linewidth}{0.9pt}\par\smallskip
  \noindent\textbf{#1.}\enspace\ignorespaces
}{%
  \par\smallskip\noindent\rule{\linewidth}{0.9pt}%
  \end{list}%
  \medskip
}

% ---------------------------------------------------------------------
%  Labelled section scheme (AC-012)
%
%  Every top-level section carries a role tag that survives into the table
%  of contents and the PDF bookmarks, so a reader can locate the new result
%  without reading linearly.  Roles:
%    exposition -- background a reader may already know and may skip
%    evidence   -- computation and measurement supporting the claim
%    contribution -- what this paper adds
%    limitations  -- what the claim does not cover
% ---------------------------------------------------------------------
\newcommand{\roletag}[1]{%
  \texorpdfstring{\ {\normalfont\small\textsc{[#1]}}}{ [#1]}%
}
\newcommand{\EXPOSITION}{\roletag{exposition}}
\newcommand{\EVIDENCE}{\roletag{evidence}}
\newcommand{\CONTRIBUTION}{\roletag{contribution}}
\newcommand{\LIMITATIONS}{\roletag{limitations}}

% ---------------------------------------------------------------------
%  Glossary machinery (AC-003)
%
%  \defterm{key}{printed term}
%      Marks the defining occurrence of a term of art in the running text.
%      It sets the term in bold italic and plants a label.
%
%  \gentry{key}{Headword}{plain-language definition}
%      A glossary line.  It back-references the section and page where
%      \defterm{key}{...} was placed.
%
%  INVARIANT: every \gentry key should have a matching \defterm in a file
%  that the paper actually \inputs.  If it does not, \gentry degrades to
%  "Definition pending" and logs a LaTeX warning rather than emitting ??,
%  so a missing anchor cannot silently break AC-001 for the other paper.
%  Grep the build log for "glossary anchor" to find any that are pending.
% ---------------------------------------------------------------------
\newcommand{\defterm}[2]{\textbf{\textit{#2}}\label{term:#1}}

\makeatletter
\newcommand{\gentry}[3]{%
  \item[#2]\ #3\ %
  \@ifundefined{r@term:#1}{%
    \@latex@warning{Missing glossary anchor for term '#1'%
      \MessageBreak (no \string\defterm{#1}{...} in this document)}%
    \emph{(Definition pending.)}%
  }{%
    \emph{(Defined in Section~\ref{term:#1}, page~\pageref{term:#1}.)}%
  }%
}
\makeatother

% ---------------------------------------------------------------------
%  Small conveniences used across both papers
% ---------------------------------------------------------------------
\newcommand{\Wmat}{\mathbf{W}}
\newcommand{\Gistar}{G_i^{*}}
\newcommand{\MoranI}{I}
\newcommand{\Ex}{\mathbb{E}}
\newcommand{\Var}{\operatorname{Var}}

% Author metadata shared by both papers.
\newcommand{\paperauthor}{Matthew Scott Tan}


\title{Selective Inference for Weights-Matrix Selection\\[0.3em]
  \large Choosing the Neighbour Parameter by Inspecting the Result Inflates the
  False-Positive Rate, and the Correction Is Derivable}
\author{\paperauthor}
\date{\today}

\hypersetup{
  pdftitle={Selective Inference for Weights-Matrix Selection},
  pdfauthor={Matthew Scott Tan},
  pdfsubject={Spatial statistics}
}

\begin{document}

\maketitle

% ---------------------------------------------------------------------
\begin{abstract}
\noindent
\textbf{[Placeholder abstract --- T-009.]}
A permutation test for spatial association builds its null distribution by
reshuffling the observed values across locations. That construction assumes the
neighbour structure was fixed before the data were seen. In practice the
neighbour parameter is routinely chosen by trying several values and keeping the
one that gives the strongest signal. This paper shows that the practice inflates
the false-positive rate by roughly a factor of three, measures the inflation as a
surface over selection rules and candidate grids rather than as a point, and
derives the correction. Numbers, figures and citations in this draft are
placeholders until the evidence tasks land; every one of them will be traceable
through \texttt{papers/CLAIMS.md}.
\end{abstract}

% ---------------------------------------------------------------------
\tableofcontents
\bigskip

% ---------------------------------------------------------------------
\section{How to read this paper}
\label{sec:roadmap}

This paper is written as lecture notes that happen to contain a new result. It
assumes a numerate reader --- algebra and reading a graph --- and no statistics
background at all. Every term of art is defined before it is used, and collected
in the glossary in Section~\ref{sec:glossary}.

Because the exposition is long by design, each section is tagged with its role,
and the tag appears in the table of contents and in the PDF bookmarks. A reader
who already knows the background can jump straight to the sections tagged
\textsc{[contribution]}.

\begin{table}[htbp]
  \centering
  \begin{tabular}{@{}llp{0.46\linewidth}@{}}
    \toprule
    Section & Role & What it is for \\
    \midrule
    \ref{sec:intro}      & exposition   & Why choosing $k$ by looking at the answer is a problem. \\
    \ref{sec:primer}     & exposition   & Everything statistical the argument needs, from scratch. \\
    \ref{sec:problem}    & exposition   & The failure, stated as a concrete procedure. \\
    \ref{sec:mechanism}  & contribution & Why selection breaks exchangeability. \\
    \ref{sec:evidence}   & evidence     & The inflation surface over rules and grids. \\
    \ref{sec:correction} & contribution & The corrected test and its derivation. \\
    \ref{sec:new}        & contribution & What is new here, against prior work. \\
    \ref{sec:limits}     & limitations  & What the claim does not cover. \\
    \bottomrule
  \end{tabular}
  \caption{Role of each section. Sections tagged \emph{contribution} carry the
    new result; the rest can be skipped by a reader who already has the
    background.}
  \label{tab:roadmap}
\end{table}

% ---------------------------------------------------------------------
\section{Introduction\EXPOSITION}
\label{sec:intro}

\textbf{[Placeholder --- T-009.]} The permutation test is the workhorse of
applied spatial statistics precisely because it makes so few assumptions. This
paper is about the one assumption it does make and that practitioners routinely
break without noticing. The statistic under test is Moran's $\MoranI$
\citep{moran1950}; the local framework is due to \citet{anselin1995}; all
sources are annotated in Section~\ref{sec:annotbib}.

Figure~\ref{fig:p2-placeholder} shows the shape of the effect.

\begin{figure}[htbp]
  \centering
  \includegraphics[width=0.82\linewidth]{placeholder.pdf}
  \caption{\textbf{[Placeholder figure.]} Schematic of the realised
    false-positive rate against the nominal five per cent, as a function of how
    many candidate $k$ values the selection rule is offered. Generated by
    \texttt{papers/shared/mkfigs.py}; replaced by the real measured surface in
    the evidence tasks.}
  \label{fig:p2-placeholder}
\end{figure}

% ---------------------------------------------------------------------
\section{A primer, from first principles\EXPOSITION}
\label{sec:primer}

This section is shared verbatim with the companion paper on component structure,
so that each paper can be read on its own. Readers who have read the other paper
can skip it.

% =====================================================================
%  papers/shared/primer.tex
%
%  PLACEHOLDER created by T-005 so the build scaffold compiles end to end.
%  T-004 OWNS THIS FILE and will replace it wholesale.  Nothing here is
%  intended as final prose.
%
%  Contract this stub establishes, which T-004 should preserve:
%    * The file is \input inside a \section, so it starts at \subsection
%      level.  It must not issue \section or \chapter.
%    * It must not assume either paper's specific context (D-006).
%    * Every term of art gets \defterm{key}{printed term} at its defining
%      occurrence, and a matching \gentry line in shared/glossary.tex.
% =====================================================================

\subsection{What this primer assumes}
\label{sec:primer-assumes}

This primer assumes a reader who is comfortable with algebra and with reading a
graph, and who has no background in statistics whatsoever. Every term of art is
defined before it is used. Terms are set in \textbf{\textit{bold italic}} at the
point of definition and collected in the glossary at the end of the paper, each
with the section number where its definition lives.

\subsection{Locations, values, and neighbours}
\label{sec:primer-weights}

Suppose we have measured one number at each of $n$ places --- a rate, a count, a
score. Write $y_i$ for the value at place $i$. Spatial statistics asks whether
the arrangement of those numbers across the map carries information that the
numbers alone do not.

To ask that question we must first say which places count as neighbours of which
others. That bookkeeping is done by a \defterm{spatial-weights-matrix}{spatial
weights matrix}: a square table $\Wmat$ with one row and one column per place,
whose entry $w_{ij}$ records how strongly place $j$ counts as a neighbour of
place $i$. A common choice sets $w_{ij}=1$ when $j$ is one of the $k$ nearest
places to $i$ and $w_{ij}=0$ otherwise.

Reading the table as a map of who-is-connected-to-whom, a group of places that
can all reach one another by hopping along nonzero weights, and which no weight
connects to anything outside the group, is a \defterm{connected-component}{connected
component} of the neighbour graph. A neighbour graph with more than one component
has been cut into pieces that never exchange information.

\subsection{Evidence against a hypothesis}
\label{sec:primer-pvalue}

Statistical tests work by contradiction. We name a dull explanation for what we
observed, then ask how surprising the observation would be if that dull
explanation were true. The dull explanation is the
\defterm{null-hypothesis}{null hypothesis}; here it is usually ``the values were
scattered across the map with no regard to location.''

The measure of surprise is the \defterm{p-value}{p-value}: the probability of
seeing a result at least as extreme as the one observed, computed under the
assumption that the null hypothesis holds. A small p-value means the observation
sits far out in the tail of what the dull explanation predicts. It does not mean
the dull explanation is false, and it is not the probability that the dull
explanation is false.

\begin{keyinsight}[Placeholder]
This is placeholder primer content supplied by the build scaffold (T-005) so
that the document compiles, the glossary resolves, and the cross-reference
machinery can be checked. T-004 replaces this file with the real
from-first-principles exposition required by AC-003.
\end{keyinsight}


% ---------------------------------------------------------------------
\section{The problem\EXPOSITION}
\label{sec:problem}

\textbf{[Placeholder --- T-009.]} The failure stated as a procedure a
practitioner would actually run: build the weights matrix at several values of
$k$, compute the statistic at each, report the best one and its p-value.

% ---------------------------------------------------------------------
\section{The mechanism\CONTRIBUTION}
\label{sec:mechanism}

\textbf{[Placeholder --- T-009.]} A permutation test builds its null by
reshuffling observed values across locations, so it assumes exchangeability ---
and exchangeability quietly requires that the weights matrix was fixed in
advance. Selecting $k$ from the data conditions the reported statistic on the
event that it was the maximum, and the unconditional null no longer describes it.

\begin{keyinsight}
The reported p-value answers the question ``how unusual is this statistic?'' when
the question actually asked was ``how unusual is the largest of several
statistics?'' The candidate statistics are strongly correlated, so the inflation
is smaller than a naive multiple-comparisons bound but far from negligible.
\end{keyinsight}

% ---------------------------------------------------------------------
\section{Evidence\EVIDENCE}
\label{sec:evidence}

\textbf{[Placeholder --- T-003 supplies the numbers, T-009 the prose.]}
Type I error inflation across at least three selection rules --- arg-max of
$\MoranI$, minimum p-value, first significant --- and at least two candidate
grids, reported as a surface rather than a point (AC-007), at a replicate count
sized for a Monte Carlo standard error below half a percentage point (A-4).

% ---------------------------------------------------------------------
\section{The correction\CONTRIBUTION}
\label{sec:correction}

\textbf{[Placeholder --- T-009.]} The corrected test, derived rather than
asserted, together with what it costs in power and when the naive test is safe
to keep using.

% ---------------------------------------------------------------------
\section{What is new here\CONTRIBUTION}
\label{sec:new}

\textbf{[Placeholder --- T-009, informed by the prior-art check in T-001.]}
This section states the novel contribution and distinguishes it from the general
post-selection inference machinery it specialises (AC-008). Where the literature
check finds work that partially preempts the claim, the claim is narrowed here
rather than the source omitted (D-008).

\begin{keyinsight}[What is new]
Placeholder. The claim stated here must survive the prior-art check, and must be
narrowed if it does not.
\end{keyinsight}

% ---------------------------------------------------------------------
\section{Limitations\LIMITATIONS}
\label{sec:limits}

\textbf{[Placeholder --- T-009.]} Which claims rest on a single dataset, which
are simulation-based, and what observation would falsify them (AC-010). In
particular, the inflation figure is measured under a specific null and a specific
family of candidate grids, and does not transfer automatically to other
selection procedures.

% ---------------------------------------------------------------------
\section{Annotated bibliography}
\label{sec:annotbib}

% =====================================================================
%  papers/shared/annotated_bibliography.tex
%
%  STUB created by T-005 so the scaffold compiles.  T-001 OWNS THIS FILE
%  and will replace it with the full annotated corpus required by AC-004.
%
%  Contract this stub establishes, which T-001 should preserve:
%    * The file is \input inside a \section, so it starts at \subsection
%      level and must not issue \section.
%    * Each source gets its own \subsection* or description item stating
%      (a) what the source does, (b) why it is credible, (c) what this
%      paper takes from it.
%    * Sources that could not be retrieved are flagged \textbf{Unverified:}.
%    * Every key annotated here must exist in shared/refs.bib and be cited
%      somewhere in the text, or BibTeX will not emit it into the
%      reference list.
% =====================================================================

The reference list that follows the annotations is generated by \textsc{BibTeX}
from \texttt{papers/shared/refs.bib}. Each entry below states what the source
does, why it is credible, and precisely what the present paper takes from it.

\begin{description}\setlength{\itemsep}{0.6em}

\item[\citet{moran1950}] \emph{What it does:} introduces the autocorrelation
  coefficient now universally called Moran's $\MoranI$, as a way of detecting
  whether a quantity measured over a region varies smoothly or erratically in
  space. \emph{Why it is credible:} the founding paper for the statistic, in
  \emph{Biometrika}, and the definition in use today is still Moran's.
  \emph{What we take:} the definition of the statistic and its ratio form.

\item[\citet{getisord1992}] \emph{What it does:} defines the family of distance
  statistics $G_i$ and $\Gistar$ for detecting local concentrations of high or
  low values. \emph{Why it is credible:} the originating paper in
  \emph{Geographical Analysis}; the statistic is implemented under this
  definition in every major spatial toolkit. \emph{What we take:} the definition
  of $\Gistar$, including the fact that it standardises against the
  \emph{global} mean and standard deviation.

\item[\citet{anselin1995}] \emph{What it does:} sets out the LISA framework,
  which decomposes a global measure of spatial association into per-location
  contributions. \emph{Why it is credible:} the standard reference for local
  spatial statistics, and the basis of the local Moran implementations in
  \texttt{spdep}, \texttt{esda} and ArcGIS. \emph{What we take:} the framing of
  local statistics as decompositions of a global one.

\end{description}

\begin{keyinsight}[Placeholder]
This is a stub annotated bibliography supplied by the build scaffold (T-005) to
prove the \textsc{BibTeX} wiring. T-001 replaces this file, and
\texttt{refs.bib}, with the retrieved and verified corpus required by AC-004 and
AC-005.
\end{keyinsight}


\addcontentsline{toc}{section}{References}
\bibliography{../shared/refs}

% ---------------------------------------------------------------------
\section{Glossary}
\label{sec:glossary}

Every term of art used in this paper, with the section in which it is defined.

% =====================================================================
%  papers/shared/glossary.tex
%
%  Glossary body shared by both papers.  The MECHANISM (\defterm, \gentry)
%  lives in shared/preamble.tex; this file holds only the entries.
%
%  Created by T-005 (scaffold).  T-004 owns the CONTENT and should extend
%  the block marked below as it writes the primer.  Keep entries in
%  alphabetical order by headword.
%
%  RULE: each \gentry{key}{Headword}{definition} needs a matching
%  \defterm{key}{...} somewhere the paper actually \inputs -- normally
%  shared/primer.tex.  If the anchor is missing the line still typesets, as
%  "Definition pending", and the build log records a warning.  Grep the log
%  for "Missing glossary anchor" before declaring a paper done.
%
%  This file is \input inside a \section, so it must not issue \section.
% =====================================================================

\begin{description}\setlength{\itemsep}{0.4em}

% ---- BEGIN entries owned by T-004 (primer terms, AC-003) --------------

\gentry{connected-component}{Connected component}{%
  A group of places in the neighbour graph that can all reach one another by
  hopping from neighbour to neighbour, and which is joined to nothing outside
  the group. A graph with two components has been cut in two.}

\gentry{null-hypothesis}{Null hypothesis}{%
  The deliberately dull explanation a statistical test tries to contradict ---
  here, that the observed values were scattered across the map with no regard
  to location.}

\gentry{p-value}{p-value}{%
  The probability of seeing a result at least as extreme as the one observed,
  computed on the assumption that the null hypothesis is true. A small p-value
  says the observation is unusual under the dull explanation; it is not the
  probability that the dull explanation is wrong.}

\gentry{spatial-weights-matrix}{Spatial weights matrix}{%
  A square table with one row and one column per place, whose entry in row $i$
  and column $j$ records how strongly place $j$ counts as a neighbour of place
  $i$. It is the object that turns a map into a graph.}

% ---- END entries owned by T-004 --------------------------------------

\end{description}


\end{document}
